% Generated from locale/tr/content/lambda-calculus/lambda-definability/partial-recursive-functions.tex; do not hand-edit.


\olfileid[tr]{lam}{ldf}{par}
\olsection{Kısmi Özyinelemeli Fonksiyonlar \usetoken{S}{lambda definable}}

Kısmi özyinelemeli fonksiyonlar, temel fonksiyonlardan bileşke, ilkel
özyineleme ve sınırsız minimizasyonla elde edilen fonksiyonlardır. Genel
özyinelemeli fonksiyonlardan farkları, sınırsız aramada kullanılan
fonksiyonların düzenli olmasının gerekmemesidir. Düzenlilik istenmediğinde
minimizasyonla tanımlanan fonksiyonlar bazı durumlarda tanımsız olabilir.

İlk bakışta bütün tümel genel özyinelemeli fonksiyonların !!{lambda definable}
olduğunu göstermek için kullanılan yöntemlerle bütün kısmi özyinelemeli
fonksiyonların da !!{lambda definable} olduğu ispatlanabilir gibi görünür.
Örneğin $F,G$ terimleri sırasıyla $f,g$'yi !!{lambda define}.
Bu durumda $\lambd[x][F(Gx)]$ terimi $f$ ile $g$'nin bileşkesini
!!{lambda define}. Ancak fonksiyonlar kısmi olduğunda sorun çıkar.
$g(x)$ tanımsızsa $Gx$'in normal biçimi yoktur. Çoğu durumda $F(Gx)$'in de
normal biçimi olmaması, istediğimiz sonuçtur. Fakat $F$'nin
$\lambd[x][\lambd[y][y]]$ olduğunu düşünün; bu durumda $F(Gx)$'in normal
biçimi $\lambd[y][y]$ vardır.

Bu sorun aşılmaz değildir. Bütün kısmi özyinelemeli fonksiyonları, tanımsız
değerleri normal biçimi olmayan terimlerle gösterecek biçimde
!!{lambda define} eden yöntemler vardır. Ancak bu yöntemler, genel
özyinelemeli fonksiyonlar için kullandığımız yaklaşımdan daha karmaşık ve daha
az sezgiseldir. Teoremi burada ispatsız kaydediyoruz:

\begin{thm}
  Bütün kısmi özyinelemeli fonksiyonlar !!{lambda definable}.
\end{thm}

